\documentclass[UTF8,zihao=-4,oneside]{ctexrep}

\usepackage[a4paper,top=25mm,bottom=25mm,left=28mm,right=24mm,headheight=15pt]{geometry}
\usepackage{fontspec}
\usepackage{xeCJK}
\usepackage{unicode-math}
\usepackage{xcolor}
\usepackage{setspace}
\usepackage{indentfirst}
\usepackage{titlesec}
\usepackage{fancyhdr}
\usepackage{graphicx}
\usepackage{float}
\usepackage{caption}
\usepackage{subcaption}
\usepackage{booktabs}
\usepackage{longtable}
\usepackage{array}
\usepackage{multirow}
\usepackage{tabularx}
\usepackage{siunitx}
\usepackage{amsmath}
\usepackage{amssymb}
\usepackage{bm}
\usepackage{enumitem}
\usepackage{tikz}
\usepackage{hyperref}
\usepackage[nameinlink,noabbrev]{cleveref}
\usepackage[
  backend=biber,
  style=gb7714-2015,
  sorting=none
]{biblatex}

\addbibresource{references.bib}

\IfFontExistsTF{Times New Roman}
  {\setmainfont{Times New Roman}}
  {\setmainfont{TeX Gyre Termes}}
\IfFontExistsTF{SimSun}
  {\setCJKmainfont{SimSun}}
  {\setCJKmainfont{FandolSong-Regular}}
\IfFontExistsTF{SimHei}
  {\setCJKsansfont{SimHei}}
  {\setCJKsansfont{FandolHei-Regular}}
\IfFontExistsTF{XITS Math}
  {\setmathfont{XITS Math}}
  {\setmathfont{Latin Modern Math}}

\definecolor{documentblack}{RGB}{0,0,0}
\color{documentblack}
\hypersetup{
  colorlinks=true,
  linkcolor=documentblack,
  citecolor=documentblack,
  urlcolor=documentblack,
  pdfborder={0 0 0},
  pdftitle={岩土工程项目技术报告}
}

\setlength{\parindent}{2em}
\setlength{\parskip}{0pt}
\setstretch{1.5}
\setlist{nosep,leftmargin=2em}
\sisetup{
  detect-all,
  per-mode=symbol,
  range-phrase={\ensuremath{\sim}},
  separate-uncertainty=true
}
\captionsetup{
  font=small,
  labelfont=bf,
  labelsep=quad,
  textfont=normalfont
}
\renewcommand{\thefigure}{\thechapter-\arabic{figure}}
\renewcommand{\thetable}{\thechapter-\arabic{table}}
\renewcommand{\theequation}{\thechapter-\arabic{equation}}

\titleformat{\chapter}[hang]
  {\centering\bfseries\zihao{3}}
  {第\chinese{chapter}章}{1em}{}
\titleformat{\section}
  {\bfseries\zihao{4}}
  {\thesection}{0.8em}{}
\titleformat{\subsection}
  {\bfseries\zihao{-4}}
  {\thesubsection}{0.8em}{}
\titlespacing*{\chapter}{0pt}{-12pt}{18pt}
\titlespacing*{\section}{0pt}{12pt}{6pt}
\titlespacing*{\subsection}{0pt}{8pt}{4pt}

\pagestyle{fancy}
\fancyhf{}
\fancyhead[C]{\zihao{5} 岩土工程项目技术报告}
\fancyfoot[C]{\thepage}
\renewcommand{\headrulewidth}{0.4pt}
\renewcommand{\footrulewidth}{0pt}

\newcommand{\ProjectName}{××项目}
\newcommand{\ReportName}{岩土工程项目技术报告}
\newcommand{\ClientName}{××建设单位}
\newcommand{\DesignerName}{××设计研究院有限公司}
\newcommand{\ProjectStage}{施工图设计}
\newcommand{\ProjectNumber}{GT-2026-001}
\newcommand{\ReportDate}{2026年9月}

\newcommand{\placeholder}[1]{\textcolor{documentblack}{【#1】}}
\newcommand{\unit}[1]{\,\si{#1}}
\newcommand{\safetyfactor}{K_{\mathrm{s}}}

\begin{document}

\begin{titlepage}
  \thispagestyle{empty}
  \centering
  \vspace*{18mm}
  {\zihao{2}\bfseries \ProjectName\par}
  \vspace{12mm}
  {\zihao{1}\bfseries \ReportName\par}
  \vspace{8mm}
  {\zihao{3}（\ProjectStage）\par}
  \vfill
  \renewcommand{\arraystretch}{1.6}
  \begin{tabular}{>{\raggedleft\arraybackslash}p{36mm}p{80mm}}
    项目编号： & \ProjectNumber \\
    建设单位： & \ClientName \\
    编制单位： & \DesignerName \\
    编制日期： & \ReportDate \\
  \end{tabular}
  \vspace*{18mm}
\end{titlepage}

\chapter*{编制人员}
\addcontentsline{toc}{chapter}{编制人员}
\begin{table}[H]
  \centering
  \renewcommand{\arraystretch}{1.7}
  \begin{tabularx}{\textwidth}{>{\centering\arraybackslash}p{25mm}X>{\centering\arraybackslash}p{35mm}}
    \toprule
    岗位 & 姓名 & 签字 \\
    \midrule
    审定 & \placeholder{填写} & \\
    审核 & \placeholder{填写} & \\
    项目负责人 & \placeholder{填写} & \\
    编制 & \placeholder{填写} & \\
    \bottomrule
  \end{tabularx}
\end{table}

\chapter*{报告说明}
\addcontentsline{toc}{chapter}{报告说明}
本报告适用于岩土工程项目设计成果编制。正式提交前应核对勘察资料、设计边界、规范版本、计算模型、图纸编号及各专业接口。所有参数均应注明来源，经验值应说明采用依据与适用条件。

\clearpage
\pagenumbering{Roman}
\tableofcontents
\clearpage
\listoffigures
\addcontentsline{toc}{chapter}{插图目录}
\clearpage
\listoftables
\addcontentsline{toc}{chapter}{表格目录}
\clearpage
\pagenumbering{arabic}

\chapter{工程概况}

\section{项目概况}
\ProjectName 位于\placeholder{项目地点}。拟建内容包括\placeholder{建筑物、构筑物、道路及场地工程}。本次工作范围包括\placeholder{边坡支护、基坑支护、场地平整、地基处理等}。

\section{任务范围}
本次设计主要完成以下工作：
\begin{enumerate}
  \item 分析场地工程地质与水文地质条件；
  \item 确定岩土设计参数及设计工况；
  \item 开展方案比选、稳定性及承载力验算；
  \item 提出构造、施工、检测、监测及风险控制要求。
\end{enumerate}

\section{设计条件与边界}
\begin{table}[H]
  \centering
  \caption{主要设计条件}
  \label{tab:design-condition}
  \begin{tabularx}{\textwidth}{p{35mm}Xp{35mm}}
    \toprule
    项目 & 取值或说明 & 来源 \\
    \midrule
    设计使用年限 & \placeholder{填写} & \placeholder{任务书或规范} \\
    工程安全等级 & \placeholder{填写} & \placeholder{判定依据} \\
    抗震设防烈度 & \placeholder{填写} & \placeholder{地勘或规范} \\
    地下水控制水位 & \placeholder{填写} & \placeholder{地勘或专项资料} \\
    周边荷载 & \placeholder{填写} & \placeholder{总图或结构资料} \\
    \bottomrule
  \end{tabularx}
\end{table}

\chapter{编制依据}

\section{基础资料}
\begin{enumerate}
  \item \placeholder{岩土工程勘察报告名称、编号、日期及版本}；
  \item \placeholder{总平面图、建筑图、结构图及竖向设计资料}；
  \item \placeholder{测量、监测、试验及现场调查资料}；
  \item 建设单位提供的其他有效资料。
\end{enumerate}

\section{采用标准}
引用标准时应写明标准全称、代号及年号。示例包括《建筑边坡工程技术规范》GB 50330—2013\cite{gb50330}、《岩土锚杆与喷射混凝土支护工程技术规范》GB 50086—2015\cite{gb50086}和《建筑基坑支护技术规程》JGJ 120—2012\cite{jgj120}。正式项目应根据任务类型核查现行版本及地方标准。

\chapter{工程地质条件}

\section{地形地貌}
\placeholder{说明场地地形、地貌单元、现状高程、坡向、沟谷及人工改造情况。}

\section{地层岩性}
\begin{longtable}{>{\centering\arraybackslash}p{16mm}p{30mm}p{24mm}p{22mm}p{22mm}p{28mm}}
  \caption{岩土层主要特征}\label{tab:strata}\\
  \toprule
  层号 & 岩土名称 & 状态或风化程度 & 层厚/m & 承载力特征值/kPa & 工程特性 \\
  \midrule
  \endfirsthead
  \multicolumn{6}{l}{续表~\thetable}\\
  \toprule
  层号 & 岩土名称 & 状态或风化程度 & 层厚/m & 承载力特征值/kPa & 工程特性 \\
  \midrule
  \endhead
  \bottomrule
  \endfoot
  ① & \placeholder{填土} & \placeholder{松散} & \placeholder{填写} & \placeholder{填写} & \placeholder{填写} \\
  ② & \placeholder{粉质黏土} & \placeholder{可塑} & \placeholder{填写} & \placeholder{填写} & \placeholder{填写} \\
  ③ & \placeholder{岩层} & \placeholder{强风化} & \placeholder{填写} & \placeholder{填写} & \placeholder{填写} \\
\end{longtable}

\section{地下水与不良地质作用}
\placeholder{说明地下水类型、稳定水位、季节变化、腐蚀性、不良地质作用和特殊性岩土。}

\chapter{设计参数与工况}

\section{设计参数}
\begin{table}[H]
  \centering
  \caption{岩土设计参数建议值}
  \label{tab:parameters}
  \resizebox{\textwidth}{!}{
  \begin{tabular}{cccccccc}
    \toprule
    岩土层 & 重度 $\gamma$/(\si{\kilo\newton\per\cubic\metre}) &
    黏聚力 $c$/kPa & 内摩擦角 $\varphi$/\si{\degree} &
    地基承载力 $f_{ak}$/kPa & 锚固体黏结强度/kPa &
    桩侧阻力/kPa & 参数来源 \\
    \midrule
    \placeholder{土层1} & \placeholder{填写} & \placeholder{填写} & \placeholder{填写} &
    \placeholder{填写} & \placeholder{填写} & \placeholder{填写} & \placeholder{地勘} \\
    \placeholder{土层2} & \placeholder{填写} & \placeholder{填写} & \placeholder{填写} &
    \placeholder{填写} & \placeholder{填写} & \placeholder{填写} & \placeholder{地勘} \\
    \bottomrule
  \end{tabular}}
\end{table}

\section{计算工况}
\begin{table}[H]
  \centering
  \caption{计算工况组合}
  \label{tab:cases}
  \begin{tabularx}{\textwidth}{cXXX}
    \toprule
    工况 & 地下水或降雨条件 & 荷载条件 & 验算项目 \\
    \midrule
    1 & 天然状态 & 永久作用与可变作用 & 稳定、承载、变形 \\
    2 & 暴雨或不利水位 & 不利荷载组合 & 稳定、抗滑、抗倾覆 \\
    3 & 地震状态 & 地震作用组合 & 整体稳定与结构承载 \\
    \bottomrule
  \end{tabularx}
\end{table}

\chapter{方案比选}

\section{备选方案}
\placeholder{列明方案构成、适用性、施工条件、工期、造价、耐久性及环境影响。}

\begin{table}[H]
  \centering
  \caption{方案综合比选}
  \label{tab:comparison}
  \begin{tabularx}{\textwidth}{p{27mm}XXX}
    \toprule
    比选因素 & 方案一 & 方案二 & 方案三 \\
    \midrule
    安全适用性 & \placeholder{填写} & \placeholder{填写} & \placeholder{填写} \\
    施工可行性 & \placeholder{填写} & \placeholder{填写} & \placeholder{填写} \\
    工期影响 & \placeholder{填写} & \placeholder{填写} & \placeholder{填写} \\
    经济性 & \placeholder{填写} & \placeholder{填写} & \placeholder{填写} \\
    风险与维护 & \placeholder{填写} & \placeholder{填写} & \placeholder{填写} \\
    \bottomrule
  \end{tabularx}
\end{table}

\section{推荐方案}
经综合比较，推荐采用\placeholder{方案名称}。推荐理由为\placeholder{填写关键控制因素及决策逻辑}。

\chapter{计算分析}

\section{计算模型}
\placeholder{说明计算软件、版本、计算方法、模型范围、边界条件、地下水处理、荷载施加及参数折减方式。}

\section{典型计算表达}
边坡整体稳定安全系数可统一记为
\begin{equation}
  \safetyfactor=\frac{\sum R_i}{\sum T_i},
  \label{eq:ks}
\end{equation}
式中，$R_i$ 为第 $i$ 条分块的抗滑作用，$T_i$ 为相应下滑作用。具体表达式应与采用的极限平衡方法一致。

挡土结构抗倾覆安全系数可写为
\begin{equation}
  K_0=\frac{\sum M_{\mathrm{res}}}{\sum M_{\mathrm{ovt}}},
  \label{eq:overturning}
\end{equation}
其中，$\sum M_{\mathrm{res}}$ 为抗倾覆力矩之和，$\sum M_{\mathrm{ovt}}$ 为倾覆力矩之和。

\section{结果汇总}
\begin{table}[H]
  \centering
  \caption{主要计算结果汇总}
  \label{tab:results}
  \begin{tabularx}{\textwidth}{p{35mm}p{28mm}p{28mm}Xp{22mm}}
    \toprule
    验算项目 & 计算值 & 限值 & 控制工况 & 结论 \\
    \midrule
    整体稳定安全系数 & \placeholder{填写} & \placeholder{填写} & \placeholder{填写} & 满足/不满足 \\
    抗滑安全系数 & \placeholder{填写} & \placeholder{填写} & \placeholder{填写} & 满足/不满足 \\
    抗倾覆安全系数 & \placeholder{填写} & \placeholder{填写} & \placeholder{填写} & 满足/不满足 \\
    地基承载力 & \placeholder{填写} & \placeholder{填写} & \placeholder{填写} & 满足/不满足 \\
    最大变形 & \placeholder{填写} & \placeholder{填写} & \placeholder{填写} & 满足/不满足 \\
    \bottomrule
  \end{tabularx}
\end{table}

\chapter{构造设计}

\section{支护结构}
\placeholder{说明结构断面、材料强度、保护层、锚固与连接、伸缩缝、排水孔、反滤层及耐久性构造。}

\section{排水系统}
\placeholder{说明坡顶截水、坡面排水、墙背排水、泄水孔、盲沟、集水井和下游可靠衔接。}

\section{典型示意图}
\begin{figure}[H]
  \centering
  \begin{tikzpicture}[x=1cm,y=1cm,line width=0.8pt]
    \draw (0,0) -- (11,0);
    \draw (2,0) -- (7,4.2) -- (11,4.2);
    \foreach \y in {0.9,1.8,2.7,3.6}{
      \pgfmathsetmacro{\x}{2+5*\y/4.2}
      \draw[->] (\x,\y) -- ++(-2.2,-0.3);
    }
    \draw[thick] (7,4.2) -- (7.8,4.2);
    \node at (5.3,2.7) {支护面};
    \node at (4.7,0.35) {坡脚排水};
    \node at (9,4.5) {坡顶截水};
  \end{tikzpicture}
  \caption{典型支护与排水关系示意}
  \label{fig:typical-support}
\end{figure}

\chapter{施工技术要求}

\section{总体顺序}
施工宜遵循分区、分层、分段原则，并结合支护体系确定开挖与支护时序。关键工序包括测量放线、排水先行、试验施工、分层开挖、随挖随支、质量检测、监测反馈和验收移交。

\section{材料与工艺}
\begin{enumerate}
  \item 进场材料应具有质量证明文件，并按规定见证取样；
  \item 隐蔽工程应在封闭前完成检查、记录和验收；
  \item 锚杆、土钉、灌注桩及注浆参数应通过设计要求和现场试验共同确定；
  \item 施工中发现地层、地下水或周边条件与勘察资料明显不符时，应暂停相关工序并及时反馈。
\end{enumerate}

\chapter{检测与监测}

\section{检测要求}
\begin{table}[H]
  \centering
  \caption{检测项目一览}
  \label{tab:test}
  \begin{tabularx}{\textwidth}{p{30mm}Xp{30mm}p{28mm}}
    \toprule
    对象 & 检测内容 & 数量或比例 & 判定依据 \\
    \midrule
    \placeholder{锚杆} & \placeholder{基本试验、验收试验} & \placeholder{填写} & \placeholder{标准条款} \\
    \placeholder{桩} & \placeholder{完整性、承载力} & \placeholder{填写} & \placeholder{标准条款} \\
    \placeholder{地基} & \placeholder{压实度、承载力} & \placeholder{填写} & \placeholder{标准条款} \\
    \bottomrule
  \end{tabularx}
\end{table}

\section{监测要求}
监测项目、点位、频率、预警值和控制值应结合工程等级、支护形式、周边环境与施工阶段确定。监测数据达到预警值或变化速率异常时，应及时复核并启动相应处置流程。

\chapter{风险控制与应急措施}

\section{主要风险}
\begin{enumerate}
  \item \placeholder{暴雨入渗、地下水上升或排水失效}；
  \item \placeholder{超挖、坡脚扰动或支护滞后}；
  \item \placeholder{成孔困难、塌孔、缩径或注浆不足}；
  \item \placeholder{周边建构筑物、道路及管线变形}。
\end{enumerate}

\section{应急原则}
出现异常位移、裂缝扩展、局部坍塌、涌水涌砂或支护构件受力异常时，应立即停止危险区域作业、设置警戒、疏散人员，并采取卸载、反压、回填、排水或临时加固等措施。后续处理应依据现场复核、监测数据与专项分析确定。

\chapter{结论与建议}

\begin{enumerate}
  \item \placeholder{概括工程地质条件和主要控制问题}；
  \item \placeholder{说明推荐方案与关键设计指标}；
  \item \placeholder{说明计算结果及安全储备}；
  \item \placeholder{提出施工、检测、监测和运行维护建议}。
\end{enumerate}

\printbibliography[heading=bibintoc,title={参考文献}]

\appendix
\chapter{计算书目录}
\placeholder{列出计算模型、软件版本、输入参数、工况、主要结果及完整计算书文件编号。}

\chapter{图纸与资料清单}
\begin{longtable}{cp{45mm}p{35mm}p{35mm}}
  \caption{成果文件清单}\\
  \toprule
  序号 & 文件名称 & 编号 & 版本或日期 \\
  \midrule
  \endfirsthead
  \toprule
  序号 & 文件名称 & 编号 & 版本或日期 \\
  \midrule
  \endhead
  \bottomrule
  \endfoot
  1 & \placeholder{设计图纸} & \placeholder{填写} & \placeholder{填写} \\
  2 & \placeholder{计算书} & \placeholder{填写} & \placeholder{填写} \\
  3 & \placeholder{相关专题报告} & \placeholder{填写} & \placeholder{填写} \\
\end{longtable}

\end{document}
